\chapter{Knowledge Graph Construction}
\label{ch:graph}
\section{Construction pipeline}
Figure~\ref{fig:pipeline} traces the flow from official records to a validated
graph. Work matching produces enriched paper records. Identity/affiliation
processing and topic classification then operate as parallel branches: missing
person evidence does not erase the official paper, and a topic label does not
resolve an author. Graph materialization combines accepted records, retains
evidence, and derives collaboration and topic-profile edges.

\begin{figure}[!htbp]\centering
\begin{tikzpicture}[font=\small,
 box/.style={draw=UETBlue,fill=PaleBlue,rounded corners=2pt,align=center,text width=6.3cm,minimum height=1.1cm,inner sep=4pt},
 wide/.style={box,text width=14cm}, arrow/.style={-{Latex[length=2mm]},semithick,color=UETBlue}]
\node[wide] (source) at (0,0) {\textbf{\lang{1. Source collection}{1. Thu thập nguồn}}\\\lang{Official AAAI/IJCAI rosters + cached OpenAlex metadata}{Danh mục AAAI/IJCAI chính thức + siêu dữ liệu OpenAlex lưu đệm}};
\node[wide] (match) at (0,-1.6) {\textbf{\lang{2. Paper matching and cleaning}{2. Đối sánh và làm sạch bản ghi bài báo}}\\\lang{Roster-derived DOI; title and byline-count checks}{DOI từ danh mục; kiểm tra tiêu đề và số tác giả}};
\draw[arrow] (source)--(match);
\node[box,minimum height=1.65cm] (identity) at (-3.8,-3.6) {\textbf{\lang{3--4. People and organizations}{3--4. Tác giả và tổ chức}}\\\lang{Compatible author IDs; ORCID conflicts; paper-specific affiliation triples}{ID tác giả phù hợp; xung đột ORCID; bộ ba liên kết tổ chức theo từng bài}};
\node[box,minimum height=1.65cm] (nlp) at (3.8,-3.6) {\textbf{\lang{5. NLP topic classification}{5. Phân loại chủ đề bằng NLP}}\\\lang{Title/abstract MiniLM embeddings; frozen taxonomy; labels or abstention}{Vector MiniLM từ tiêu đề/tóm tắt; hệ chủ đề cố định; gán nhãn hoặc từ chối}};
\draw[arrow] (match.south -| identity.north)--(identity.north);
\draw[arrow] (match.south -| nlp.north)--(nlp.north);
\node[wide] (graph) at (0,-5.6) {\textbf{\lang{6. Graph materialization}{6. Tạo các quan hệ đồ thị}}\\\lang{Five entity types; observed, inferred, curated and derived edges; evidence links}{Năm loại thực thể; cạnh quan sát, suy ra, biên soạn và dẫn xuất; liên kết bằng chứng}};
\draw[arrow] (identity.south)--(identity.south |- graph.north);
\draw[arrow] (nlp.south)--(nlp.south |- graph.north);
\node[wide] (validate) at (0,-7.25) {\textbf{\lang{7. Validation and storage}{7. Kiểm tra và lưu trữ}}\\\lang{Schema, endpoints, evidence, support counts; Neo4j; analysis and Explorer}{Lược đồ, đầu mút, bằng chứng, số bài hỗ trợ; Neo4j; phân tích và Explorer}};
\draw[arrow] (graph)--(validate);
\end{tikzpicture}
\caption{\lang{Data-to-graph workflow. Identity/affiliation handling and topic assignment are parallel dependencies. Frozen v4 retains the seven historical DOI exceptions described in Chapter~\ref{ch:data}.}{Quy trình từ dữ liệu đến đồ thị. Xử lý định danh/liên kết tổ chức và gán chủ đề là hai nhánh song song. V4 cố định vẫn giữ bảy trường hợp DOI lịch sử nêu ở Chương~\ref{ch:data}.}}
\label{fig:pipeline}\end{figure}


The graph is stored in JSON and loaded into an
isolated Neo4j Community 5.26 database \citep{neo4jDocs}. Named
projections support offline analysis; the Explorer provides read-only access to
entities, relationships and source evidence. All reported measurements use
the same graph and analysis configuration.

\section{NLP-based topic classification}
The project uses natural language processing (NLP) to connect papers to research topics. For each matched
work, the input is its title followed by its reconstructed abstract; title-only
processing is supported when an abstract is absent. The same frozen
\modelid{sentence-transformers/all-MiniLM-L6-v2} model encodes paper text and
leaf-topic descriptions \citep{reimers2019sentencebert,minilmModelCard}.
Descriptions combine the topic name, definition, included concepts and synonyms.
The pinned model revision is recorded in the repository configuration.

The classifier compares normalized vectors by cosine similarity. It accepts
the best leaf when its score reaches the development-calibrated threshold
$\tau\approx0.2709$. It can add up to two more leaves if they also exceed that
threshold and lie within $\delta=0.02$ of the best score. Otherwise it abstains.
This produces multi-label \texttt{HAS\_TOPIC} edges rather than forcing every
paper into one category. A supporting abstract sentence is selected by semantic
similarity and retained with a source locator; the title supplies evidence for
title-only input. Similarity scores are not calibrated probabilities.

The taxonomy contains four parent and 16 leaf topics. Its curated
\texttt{SUBTOPIC\_OF} links are separate from model-assigned paper labels.
Researcher and organization topic profiles are derived later from accepted
paper support. The classifier and thresholds remain frozen; no new topic
inference was run for this report.

\section{Entity resolution and affiliation mapping}
\textbf{Papers.} The official roster ID identifies the proceedings record;
DOI, title and byline checks link it to an OpenAlex work. Ambiguous candidates
stay outside accepted matches. The historical DOI caveat is described in
Chapter~\ref{ch:data}.

\textbf{Researchers.} At each byline position, a compatible official name and
OpenAlex authorship supply the researcher ID. The internal key contains the
OpenAlex author ID; a shared normalized name alone never merges people. ORCID is
retained when present, but is not mandatory. Conflicting checksum-valid ORCIDs
for one provider ID cause global withholding. More than one incompatible
position or a repeated author ID can withhold a paper's full authorship mapping.
Thus an accepted work match does not mean all its printed authors are resolved.

\textbf{Organizations.} Institutions must have OpenAlex IDs and appear on the
specific accepted authorship. Organization records retain normalized sector and
country metadata. The pipeline stores each researcher--paper--organization
observation before aggregating \texttt{AFFILIATED\_WITH}. Raw affiliation strings
remain source evidence; a current profile or a paper-wide institution list is
not copied to every author. A guard withholds suspicious identical multi-institution
sets copied across several authors. Unknown institutions or sectors remain unknown.

\section{Graph schema and relationship semantics}
The five node types are Researcher, Paper, Organization, Topic and Venue.
Table~\ref{tab:schema} summarizes the ten relationships. The four evidence
classes distinguish observed metadata, model inference, curated vocabulary and
derivation from accepted base edges.

\begin{table}[!htbp]\centering\small
\caption{\lang{The ten relationships and their support rules. Arrows specify direction.}{Mười quan hệ và quy tắc hỗ trợ. Mũi tên chỉ chiều quan hệ.}}
\label{tab:schema}
\setlength{\tabcolsep}{4pt}\renewcommand{\arraystretch}{1.12}
\begin{tabularx}{\textwidth}{>{\raggedright\arraybackslash}p{4.1cm}>{\raggedright\arraybackslash}p{3.7cm}Y}\toprule
\lang{Predicate / class}{Quan hệ / loại} & \lang{Endpoints}{Hai đầu quan hệ} & \lang{Meaning}{Ý nghĩa}\\\midrule
\texttt{AUTHORED} & \lang{Researcher $\to$ Paper}{Nhà nghiên cứu $\to$ Bài báo} & \lang{Observed: accepted byline identity.}{Quan sát: định danh tác giả được chấp nhận.}\\
\texttt{AFFILIATED\_WITH} & \lang{Researcher $\to$ Organization}{Nhà nghiên cứu $\to$ Tổ chức} & \lang{Observed: aggregate paper-specific affiliations.}{Quan sát: tổng hợp liên kết tổ chức theo từng bài.}\\
\texttt{PUBLISHED\_IN} & \lang{Paper $\to$ Venue}{Bài báo $\to$ Hội nghị} & \lang{Observed: official proceedings entry.}{Quan sát: bản ghi kỷ yếu chính thức.}\\
\texttt{CITES} & \lang{Paper $\to$ Paper}{Bài báo $\to$ Bài báo} & \lang{Observed: in-scope citing-to-cited reference.}{Quan sát: trích dẫn tới bài trong phạm vi.}\\
\texttt{HAS\_TOPIC} & \lang{Paper $\to$ Topic}{Bài báo $\to$ Chủ đề} & \lang{Inferred: one to three leaf labels, or abstention.}{Suy ra từ mô hình: một đến ba nhãn lá, hoặc không gán.}\\
\texttt{SUBTOPIC\_OF} & \lang{Topic $\to$ Topic}{Chủ đề $\to$ Chủ đề} & \lang{Curated: leaf-to-parent taxonomy.}{Biên soạn: chủ đề lá thuộc chủ đề cha.}\\
\texttt{COAUTHORED\_WITH} & \lang{Researcher $\leftrightarrow$ Researcher}{Nhà nghiên cứu $\leftrightarrow$ Nhà nghiên cứu} & \lang{Derived: shared accepted core paper.}{Dẫn xuất: cùng bài thuộc tập lõi, tác giả được chấp nhận.}\\
\texttt{COLLABORATED\_WITH} & \lang{Organization $\leftrightarrow$ Organization}{Tổ chức $\leftrightarrow$ Tổ chức} & \lang{Derived: distinct authors with aligned affiliations.}{Dẫn xuất: tác giả khác nhau có liên kết tổ chức tương ứng.}\\
\texttt{RESEARCHES\_TOPIC} & \lang{Researcher $\to$ Topic}{Nhà nghiên cứu $\to$ Chủ đề} & \lang{Derived: at least two distinct core papers.}{Dẫn xuất: ít nhất hai bài phân biệt thuộc tập lõi.}\\
\makecell[l]{\texttt{CONTRIBUTES\_TO\_}\\\texttt{TOPIC}} & \lang{Organization $\to$ Topic}{Tổ chức $\to$ Chủ đề} & \lang{Derived: at least three distinct core papers.}{Dẫn xuất: ít nhất ba bài phân biệt thuộc tập lõi.}\\\bottomrule
\end{tabularx}\end{table}

\FloatBarrier

Each semantic edge has a deterministic ID, accepted status, endpoint IDs,
evidence references and construction-method metadata. Symmetric collaboration pairs
are stored once in canonical ID order and analyzed as undirected. Source
confirmations and reverse copies do not create additional edges.

For a paper with $k$ official authors, each accepted coauthor pair receives
fractional weight $1/(k-1)$. Unresolved authors remain in $k$. For organization
collaboration, two distinct authors must support two distinct institutions on
that paper; a sole dual-affiliated author cannot create the edge. A paper
supporting $q$ eligible organization pairs contributes $1/q$ to each pair.
Both relationships also retain the distinct supporting-paper count.
\texttt{RESEARCHES\_TOPIC} requires two distinct core papers;
\texttt{CONTRIBUTES\_TO\_TOPIC} requires three. Repeated evidence for one paper
does not increase those counts.

\Needspace{18cm}
\section{Graph construction example}
Consider \emph{ECG ODE-GAN: Learning Ordinary Differential Equations of ECG
Dynamics via Generative Adversarial Learning} (AAAI 2021). Its
\href{https://ojs.aaai.org/index.php/AAAI/article/view/16086}{official record}
becomes \sourcepath{paper:aaai:2021:16086}. DOI
\href{https://doi.org/10.1609/aaai.v35i1.16086}{\texttt{10.1609/aaai.v35i1.16086}}
matches \href{https://openalex.org/W3176135934}{OpenAlex W3176135934} with exact
title agreement and three authors in both records.

Accepted authorships link the paper to Tomer Golany (A5022512211, position 1) and Kira Radinsky
(A5029708595, position 3). The official middle author, Daniel Freedman, has no
accepted authorship edge in this graph. Two accepted paper-specific observations
link Golany and Radinsky to Technion (I174306211). The classifier assigns
\emph{Generative modeling}, with stored cosine similarity 0.3753 and an abstract
excerpt about learning dynamics through a generative adversarial network.

\begin{figure}[!htbp]\centering
\begin{tikzpicture}[font=\small,
 box/.style={draw=UETBlue,fill=PaleBlue,rounded corners=2pt,align=center,text width=4.0cm,minimum height=0.9cm,inner sep=4pt},
 arrow/.style={-{Latex[length=1.8mm]},color=UETBlue},lab/.style={font=\footnotesize,fill=white,inner sep=1.5pt}]
\node[box,text width=5.5cm] (org) at (0,3.4) {Technion\\\texttt{I174306211}};
\node[box] (a) at (-4.4,1.6) {Tomer Golany\\\texttt{A5022512211}};
\node[box] (b) at (4.4,1.6) {Kira Radinsky\\\texttt{A5029708595}};
\node[box,text width=5.5cm] (paper) at (0,-0.2) {ECG ODE-GAN (AAAI 2021)\\\texttt{paper:aaai:2021:16086}};
\node[box] (venue) at (-4.4,-2.2) {AAAI\\\texttt{venue:aaai}};
\node[box] (topic) at (4.4,-2.2) {\lang{Generative modeling}{Mô hình sinh}\\{\footnotesize\texttt{topic:}\\[-2pt]\texttt{generative-modeling}}};
\draw[arrow] (a.north)--node[lab,above]{R2}(org.west);
\draw[arrow] (b.north)--node[lab,above]{R2}(org.east);
\draw[arrow] (a.south)--node[lab,below]{R1}(paper.west);
\draw[arrow] (b.south)--node[lab,above]{R1}(paper.east);
\draw[dashed,color=UETGold,thick] (a.east)--node[lab,above]{R7}(b.west);
\draw[arrow] (paper.south west)--node[lab,below]{R3}(venue.north);
\draw[arrow,dashed] (paper.south east)--node[lab,above]{R5}(topic.north);
\end{tikzpicture}
\par\smallskip
{\footnotesize\lang{Relationships}{Quan hệ}: R1 \texttt{AUTHORED}; R2 \texttt{AFFILIATED\_WITH};
R3 \texttt{PUBLISHED\_IN};\\
R5 \texttt{HAS\_TOPIC}; R7 \texttt{COAUTHORED\_WITH}.\par}
\caption{\lang{Accepted relationships for the example. IDs are shortened to provider IDs where shown. Solid arrows are observed; dashed topic arrow is inferred; the dashed undirected line is derived. Affiliation edges aggregate paper-specific observations; both displayed observations exist for this paper.}{Các quan hệ được chấp nhận trong ví dụ. Một số ID được rút gọn thành ID nhà cung cấp. Mũi tên liền là quan sát; mũi tên nét đứt tới chủ đề là suy ra; đường vô hướng nét đứt là dẫn xuất. Cạnh liên kết tổ chức tổng hợp các quan sát theo bài; cả hai quan sát được thể hiện đều tồn tại ở bài này.}}
\label{fig:example}\end{figure}

\FloatBarrier

The displayed authorship, affiliation, venue, topic and coauthor edges all exist
in the graph. This paper contributes $1/(3-1)=0.5$ to the accepted coauthor
pair's fractional weight; the edge can also aggregate other papers. Both
accepted authors point to the same institution, so this example supplies no
organization--organization collaboration. One paper alone also cannot establish
a researcher-topic or organization-topic profile edge. Full stable IDs,
accepted edge IDs and the two affiliation observations are recorded in
\sourcepath{report/latex/data/construction-example.json}.
